\section*{Hoofdstuk \ref{ch:b2:wittig} --- Wittig en andere reacties die C=C vormen}

\begin{solution}{exo:b2:wittig:1}
\ce{Ph3P + CH3CH2Br -> [Ph3PCH2CH3]+ + Br-} (bimoleculaire substitutie, verhitten in een oplosmiddel), daarna
butyllithium in droge ether of THF onder inert gas: \ce{[Ph3PCH2CH3]+ + BuLi -> Ph3P=CHCH3 + BuH + Li+}.
\end{solution}

\begin{solution}{exo:b2:wittig:2}
Het ylide \ce{Ph3P=CH2} en benzaldehyde geven fenyletheen (styreen), \ce{C6H5CH=CH2}, en
trifenylfosfineoxide.
\end{solution}

\begin{solution}{exo:b2:wittig:3}
\ce{Ph3P=CHCOOEt} en \ce{Ph3P=CHCOCH3}: gestabiliseerd. \ce{Ph3P=CHCH3}: niet-gestabiliseerd, het ylide dat
butyllithium vraagt (of natriumhydride, natriumamide). \ce{Ph3P=CHC6H5}: arylgestabiliseerd, in de praktijk gemaakt met een
alkoxide of hydroxide, en het geeft mengsels van geometrieën.
\end{solution}

\begin{solution}{exo:b2:wittig:4}
Ethyl-(\textit{E})-pent-2-enoaat, \ce{CH3CH2CH=CHCOOEt}, met natriumdiethylfosfaat.
\end{solution}

\begin{solution}{exo:b2:wittig:5}
Wittig: het ylide \ce{Ph3P=CH2}, gemaakt uit methyltrifenylfosfoniumjodide en butyllithium, geeft met cyclohexanon alleen
methyleencyclohexaan, met de dubbele binding buiten de ring. Dehydratatie van 1-methylcyclohexaan-1-ol verloopt
via een tertiair carbokation en verliest het waterstofatoom dat het meest gesubstitueerde alkeen geeft (Zaitsev):
vooral 1-methylcyclohexeen, het verkeerde isomeer.
\end{solution}

\begin{solution}{exo:b2:wittig:6}
Beide sneden zijn dezelfde: benzaldehyde en het benzylideenylide \ce{Ph3P=CHC6H5}. Dat ylide is alleen
arylgestabiliseerd en geeft een mengsel van (\textit{E}) en (\textit{Z}). Voor zuiver (\textit{E}) gebruik je de HWE-reactie
van diethylbenzylfosfonaat met benzaldehyde.
\end{solution}

\begin{solution}{exo:b2:wittig:7}
Met \ce{Ph3P=CHCH3}: (\textit{Z})-pent-2-een, \ce{CH3CH2CH=CHCH3}, vooral (\textit{Z}). Met
\ce{Ph3P=CHCOOEt}: ethyl-(\textit{E})-pent-2-enoaat, vooral (\textit{E}).
\end{solution}

\begin{solution}{exo:b2:wittig:8}
\ce{Ph3P=CH2 + C6H10O -> C7H12 + Ph3PO}. Massa's: methyleencyclohexaan \qty{96.173}{g/mol},
trifenylfosfineoxide \qty{278.291}{g/mol}; atoomeconomie $96.173/(96.173 + 278.291) =
96.173/374.464 \approx 25.7~\%$. Drie kwart van de massa van de reagentia vertrekt als fosfineoxide.
% ledger: aw:C, aw:H, aw:O, aw:P
\end{solution}

\begin{solution}{exo:b2:wittig:9}
Fosfor, nucleofiel, verdringt bromide:
\begin{equation*}
\ce{P(OEt)3 + BrCH2COOEt -> [(EtO)3PCH2COOEt]+ + Br-}.
\end{equation*}
Daarna
valt bromide een koolstofatoom van een ethylgroep van het fosfoniumion aan (bimoleculaire substitutie), wat
broomethaan vrijmaakt en de binding \ce{P=O} vormt:
\begin{equation*}
\ce{[(EtO)3PCH2COOEt]+ + Br- -> (EtO)2P(O)CH2COOEt + EtBr}.
\end{equation*}
\end{solution}

\begin{solution}{exo:b2:wittig:10}
\ce{OHC-CH=CH-CHO + 2Ph3P=CHC6H5 -> C6H5CH=CH-CH=CH-CH=CHC6H5 + 2Ph3PO}: twee equivalenten van het
benzylideenylide (uit benzyltrifenylfosfoniumchloride en een base), één voor elke aldehydegroep.
De nieuwe dubbele bindingen ontstaan als mengsels van geometrieën; het all-(\textit{E})-isomeer, het
stabielste, wordt geïsoleerd door kristallisatie, of het mengsel wordt geïsomeriseerd.
\end{solution}

\begin{solution}{exo:b2:wittig:11}
\omterm{def:b2:enolates-aldol:aldol}{Aldol}: benzaldehyde met propanon in overmaat en natriumhydroxide (\cref{ch:b2:enolates-aldol});
goedkope reagentia, water als bijproduct, (\textit{E})-product, maar de condensatie kan twee keer verlopen
(dibenzalaceton) als propanon niet in overmaat is. Wittig: benzaldehyde met het \omterm{def:b2:wittig:ylide}{gestabiliseerde ylide}
\ce{Ph3P=CHCOCH3}; (\textit{E}), één product, geen dubbele condensatie, maar het ylide is duur en
trifenylfosfineoxide moet verwijderd worden. De aldolroute is hier de betere.
\end{solution}

\begin{solution}{exo:b2:wittig:12}
Wittig: butanal met \ce{Ph3P=CHCH2CH2CH3} (uit 1-broombutaan, daarna butyllithium), zoutvrij:
vooral (\textit{Z}); bijproduct trifenylfosfineoxide, en wat (\textit{E})-isomeer. Alkyn:
oct-4-yn (uit ethyn door twee alkyleringen met 1-broompropaan en natriumamide), gehydrogeneerd over een
vergiftigde palladiumkatalysator (\cref{prop:b2:alkene-redox:syn-hydrogenation}): alleen (\textit{Z}); de
bijproducten zijn natriumbromide en ammoniak uit de alkyleringen.
\end{solution}

\begin{solution}{pb:b2:wittig:1}
\textbf{1.} \ce{C23H46}, dus \ce{C_nH_{2n}} met $n = 23$: één dubbele binding (of één ring), hier een
\ce{C=C}.
\textbf{2.} De dubbele binding \ce{C9=C10}.
\textbf{3.} Nonanal \ce{CH3(CH2)7CHO} met het ylide uit een 1-halogeentetradecaan; of tetradecanal
\ce{CH3(CH2)12CHO} met het ylide uit een 1-halogeennonaan.
\textbf{4.} 1-Broomtetradecaan, \ce{CH3(CH2)13Br}.
\textbf{5.} Bimoleculaire substitutie door het omvangrijke trifenylfosfine verloopt snel op een ongehinderd primair
koolstofatoom en concurreert niet met eliminatie.
\textbf{6.} \ce{Ph3P + C14H29Br -> [Ph3PC14H29]+ + Br-}, tetradecyltrifenylfosfoniumbromide.
\textbf{7.} Butyllithium (of natriumhydride, natriumamide); hydroxide is een te zwakke base voor een
niet-gestabiliseerd ylide, en het water dat het meebrengt, zou het ylide protoneren.
\textbf{8.} Nee, het draagt alleen een alkylgroep: vooral (\textit{Z}) (\cref{prop:b2:wittig:stereo}).
\textbf{9.} Een vierring \ce{P-C-C-O}: fosfor gebonden aan het koolstofatoom dat de keten \ce{C13H27} draagt,
dat koolstofatoom aan het koolstofatoom dat de keten \ce{C8H17} draagt, en dat aan zuurstof, dat aan fosfor
gebonden is.
\textbf{10.} Butyllithium en het ylide reageren met water en zuurstof.
\textbf{11.} De dubbele binding verbindt het vroegere carbonylkoolstofatoom van nonanal (C9 van het product) met het
vroegere ylidekoolstofatoom (C10), en geen enkele andere binding verschuift (\cref{prop:b2:wittig:regiospecific}).
\textbf{12.} Een carbokation op C9, dat een waterstofatoom van C8 of C10 kan verliezen en kan omleggen: een mengsel van
tricos-8-een en tricos-9-een, elk als (\textit{E}) en (\textit{Z}), vooral (\textit{E}).
\textbf{13.} \ce{Ph3P=CHC13H27 + C8H17CHO -> C8H17CH=CHC13H27 + Ph3PO}.
\textbf{14.} Doel \ce{C23H46}: \qty{322.621}{g/mol}; nonanal \ce{C9H18O}: \qty{142.242}{g/mol}; zout
\ce{C32H44BrP}: \qty{539.582}{g/mol}; \ce{Ph3PO}, \ce{C18H15OP}: \qty{278.291}{g/mol}.
\textbf{15.} Het ylide \ce{C32H43P} weegt $539.582 - 80.912 = \qty{458.670}{g/mol}$; met nonanal is dat
\qty{600.912}{g/mol} reagentia voor \qty{322.621}{g/mol} product: $322.621/600.912 \approx
53.7~\%$.
\textbf{16.} Via de polariteit: de apolaire koolwaterstof gaat door een korte silicakolom met een
apolair elutiemiddel terwijl het oxide achterblijft; of het oxide kristalliseert uit een koud apolair oplosmiddel.
\textbf{17.} De zeer sterke binding \ce{P=O} die gevormd wordt, weegt zwaarder dan de verbroken binding \ce{P-C}
(\cref{prop:b2:wittig:driving}).
\textbf{18.} Ze vraagt een \omterm{def:b2:wittig:hwe}{fosfonaat} met een elektronenzuigende groep op het reagerende koolstofatoom, wat ontbreekt
in een gewone alkylketen, en ze geeft het (\textit{E})-alkeen.
\textbf{19.} Tricos-9-yn, \ce{CH3(CH2)7C#C(CH2)12CH3}, met waterstof over een vergiftigde
palladiumkatalysator (palladium op calciumcarbonaat met een loodzout en chinoline).
\textbf{20.} Dec-1-yn, \ce{CH3(CH2)7C#CH}, gedeprotoneerd door natriumamide, daarna 1-broomtridecaan
\ce{CH3(CH2)12Br}: $8 + 2 + 13 = 23$ koolstofatomen, de drievoudige binding tussen C9 en C10.
\textbf{21.} Verdere hydrogenering zou tricosaan geven, zonder dubbele binding; de vergiftigde katalysator
bindt het alkyn veel sterker dan het alkeen, dat het oppervlak verlaat vóór een tweede additie.
\textbf{22.} Beide vragen twee stappen uit beschikbare stukken. De alkynroute geeft het (\textit{Z})-isomeer
zuiver en laat alleen zouten achter; de wittigroute geeft wat (\textit{E})-isomeer en een grote massa
trifenylfosfineoxide.
\textbf{23.} 100~\%: \ce{C23H44 + H2 -> C23H46}, elk atoom komt in het product terecht.
\textbf{24.} $278.291/322.621 = \boldsymbol{\approx \qty{0.86}{g}}$ trifenylfosfineoxide per
gram feromoon.
\end{solution}
